
%%=============================================================================
%% Analyse en ontwerp
%%=============================================================================
\chapter{\IfLanguageName{dutch}{Analyse en ontwerp}{Analysis and design}}%
\label{ch:ontwerp}
Voor de ontwikkeling van de proof of concept (PoC) werd eerst onderzocht
welke behoefte aan de basis ligt van het project en welke functionaliteiten
nodig zijn om deze behoefte te beantwoorden. De aanleiding is een bestaand
backlog-item binnen Vesalius.ai rond het flexibeler beheren van automatisch
verstuurde e-mails. De functionele verwachtingen werden verder verduidelijkt
op basis van de bespreking met de CEO en de beschikbare projectafspraken.

De belangrijkste ontwerpuitdaging bestaat erin gebruikers de inhoud en opmaak
van e-mails te laten aanpassen zonder de technische structuur van de
e-mailverwerking volledig vrij te geven. Daarbij moet rekening worden gehouden
met de veilige verwerking van HTML, de vervanging van dynamische placeholders
en de beperkingen van verschillende e-mailclients.

Om deze aspecten onafhankelijk van het bestaande productieplatform te kunnen
onderzoeken, wordt gekozen voor een afzonderlijke testopstelling met een
Angular-frontend en een mock-API. Hierdoor kunnen de editor, templateopslag
en renderingpipeline afzonderlijk worden ontwikkeld en gevalideerd. De
architectuur en functionele scope zijn bewust beperkt tot de onderdelen die
nodig zijn om de technische haalbaarheid van de oplossing aan te tonen.

\section{\IfLanguageName{dutch}{Vereisten}{Requirements}}%
\label{sec:vereisten}
De functionele vereisten werden afgeleid uit de bestaande productbehoefte en
verder afgestemd tijdens de bespreking met de CEO. De PoC moet aantonen dat
gebruikers e-mailtemplates visueel kunnen aanpassen en dat de gegenereerde
e-mails correct worden verwerkt en getest, zonder daarvoor de productieomgeving
van Vesalius.ai te wijzigen.

De belangrijkste functionele vereisten zijn:
\begin{itemize}
    \item \textbf{E-mailinhoud bewerken:} gebruikers kunnen de tekst en
    opmaak van een e-mail aanpassen via een visuele editor, zonder rechtstreeks
    HTML-code te moeten schrijven.

    \item \textbf{Templatebeheer:} aangepaste templates kunnen worden
    opgeslagen en opnieuw geladen op basis van e-mailtype, taal en configuratie.

    \item \textbf{Standaardtemplates:} wanneer geen aangepaste template
    beschikbaar is, kan de testopstelling terugvallen op een standaardtemplate.

    \item \textbf{Dynamische placeholders:} ondersteunde placeholders kunnen
    in de editor worden ingevoegd en tijdens de rendering worden vervangen
    door fictieve gegevens.

    \item \textbf{Vaste e-mailschil:} de gegenereerde e-mail wordt voorzien
    van een vooraf bepaalde structuur met onder meer logo, kleuren en footer.
    De PoC richt zich voornamelijk op het aanpassen van de e-mailbody.

    \item \textbf{Rendering:} de HTML wordt opgeschoond, placeholders worden
    veilig vervangen en CSS wordt inline verwerkt. Daarnaast wordt een
    tekstuele versie van de e-mail gegenereerd.

    \item \textbf{Voorbeeldweergave en testverzending:} gebruikers kunnen
    de gerenderde e-mail bekijken en een testmail versturen naar de lokale
    testmailserver.

    \item \textbf{Ondersteunde e-mailtypes:} de PoC demonstreert de werking
    aan de hand van afspraakbevestigingen, herinneringen en annulaties.
\end{itemize}

Naast de functionele vereisten zijn er verschillende niet-functionele
vereisten. De oplossing moet onderhoudbaar en testbaar zijn en een duidelijke
scheiding behouden tussen de frontend en de mock-API. De HTML-verwerking moet
rekening houden met onveilige invoer en ongeldige placeholders. Ook de
weergave in verschillende e-mailclients vormt een belangrijk aandachtspunt,
omdat HTML en CSS niet overal op dezelfde manier worden geïnterpreteerd.

De testopstelling moet bovendien reproduceerbaar zijn en uitsluitend gebruikmaken
van fictieve gegevens. De PoC mag geen afhankelijkheid creëren van echte
patiëntgegevens of productiecredentials.

\section{\IfLanguageName{dutch}{Architectuur}{Architecture}}%
\label{sec:architectuur}
De PoC wordt ontwikkeld als een afzonderlijke applicatie die losstaat van
de bestaande Vesalius-codebase. In tegenstelling tot een rechtstreekse
uitbreiding van het productieplatform, maakt deze opstelling gebruik van
een eigen frontend, een mock-API en een lokale testmailserver.

De Angular-frontend vormt de gebruikersinterface en bevat de visuele
e-maileditor, de templateconfiguratie, de voorbeeldweergave en de bediening
voor het versturen van testmails. De frontend communiceert via een API met
de mock-backend, die instaat voor het opslaan en laden van templates en
het uitvoeren van de renderingpipeline.

Bij het verwerken van een template wordt eerst de HTML opgeschoond. Vervolgens
worden de toegestane placeholders vervangen door fictieve voorbeeldgegevens
en wordt de inhoud in de vaste e-mailschil geplaatst. De CSS wordt inline
verwerkt en er wordt een tekstuele versie gegenereerd. De resulterende e-mail
kan daarna worden weergegeven in de preview of worden verstuurd naar Mailpit.

De testomgeving wordt met Docker en Docker Compose opgezet, zodat de lokale
mailserver en de benodigde diensten reproduceerbaar kunnen worden gestart.
Geautomatiseerde tests controleren de belangrijkste onderdelen van de
verwerking. Aanvullende controles in echte e-mailclients worden gebruikt om
weergaveverschillen vast te stellen.

Een vereenvoudigde voorstelling van de beoogde architectuur is:

\begin{figure}[htbp]
    \centering
    \includegraphics[width=0.9\textwidth]{email-architectuur}
    \caption[Architectuur van de e-maileditor-PoC]{%
        Beoogde architectuur van de zelfstandige PoC. De Angular-frontend
        communiceert met de mock-API voor templatebeheer en rendering.
        Mailpit ontvangt testmails en ondersteunt de controle van de
        gegenereerde e-mails.}
    \label{fig:architectuur}
\end{figure}

De belangrijkste ontwerpkeuze is de scheiding tussen de editor en de
e-mailverwerking. De frontend is verantwoordelijk voor de interactie met de
gebruiker, terwijl de mock-API de templateverwerking en het genereren van de
e-mail uitvoert. Deze scheiding maakt het mogelijk beide onderdelen
afzonderlijk te testen en vergemakkelijkt een eventuele latere integratie
met de bestaande backend.

\section{\IfLanguageName{dutch}{Gegevensmodel}{Data model}}%
\label{sec:gegevensmodel}
De PoC heeft een beperkt gegevensmodel dat voornamelijk gericht is op het
beheren en renderen van e-mailtemplates. De gegevensstructuur moet voldoende
informatie bevatten om een template terug te vinden, de juiste configuratie
te selecteren en de inhoud met fictieve gegevens te kunnen renderen.

Een template bevat conceptueel minstens de volgende gegevens:
\begin{itemize}
    \item \textbf{Identificatie:} een unieke identificatie van de template.
    \item \textbf{Configuratie:} het bijbehorende e-mailtype, de taal en
    de configuratie waarvoor de template bedoeld is.
    \item \textbf{Inhoud:} de aanpasbare e-mailbody, met de ondersteunde
    opmaak en placeholders.
    \item \textbf{Status:} informatie die nodig is om onderscheid te maken
    tussen een standaardtemplate en een aangepaste versie.
\end{itemize}

De mockgegevens bevatten daarnaast fictieve voorbeeldgegevens voor patiënten,
afspraken, artsen en organisaties. Deze gegevens worden gebruikt om de
placeholders tijdens de rendering te vervangen en om verschillende
randgevallen te kunnen testen, zoals ontbrekende namen, lange adressen en
speciale tekens.

De opslagtechnologie wordt niet vooraf vastgelegd op basis van de
productiedatabase van Vesalius.ai. Voor de PoC wordt een tijdelijke
oplossing gekozen die eenvoudig op te zetten, reproduceerbaar en voldoende
voor de testvereisten is. De precieze gegevensstructuur en opslagkeuze
worden tijdens de implementatie verder gevalideerd.

Een belangrijk aandachtspunt is dat onbekende placeholders niet stilzwijgend
als geldige gegevens mogen worden behandeld. De renderingpipeline moet
controleren welke placeholders ondersteund worden en moet rekening houden
met ontbrekende of ongeldige waarden. Op die manier blijft de gegenereerde
e-mail voorspelbaar en kunnen fouten tijdens het testen gemakkelijker
worden opgespoord.

\section{\IfLanguageName{dutch}{Privacy, veiligheid en regelgeving}
{Privacy, security and regulation}}%
\label{sec:privacy}
Hoewel de PoC gericht is op e-mailcommunicatie binnen een medische context,
wordt de testopstelling bewust gescheiden gehouden van de productieomgeving.
Alle gebruikte patiënt-, afspraak- en organisatiegegevens zijn fictief.
Hierdoor is het niet nodig om voor de demonstratie echte medische
persoonsgegevens te verwerken.

Bij de ontwikkeling wordt rekening gehouden met de principes van de
Algemene Verordening Gegevensbescherming (AVG/GDPR), waaronder
gegevensminimalisatie en vertrouwelijkheid. De PoC verwerkt uitsluitend
de gegevens die noodzakelijk zijn om de functionaliteiten van de editor
en de renderingpipeline te demonstreren. De tijdelijke opslag bevat
testgegevens en geen echte patiëntendossiers.

Een belangrijk technisch risico is het verwerken van HTML die door een
gebruiker kan worden aangepast of geplakt vanuit een andere toepassing.
HTML mag daarom niet zonder controle worden overgenomen en gerenderd.
De renderingpipeline voorziet in HTML-opschoning op basis van toegestane
elementen en attributen. Daarnaast moeten placeholders volgens een
whitelist worden verwerkt, waarbij waarden veilig worden geëncodeerd
om ongewenste HTML-injectie te voorkomen. Deze maatregelen verkleinen
het risico op cross-site scripting (XSS) en andere problemen door
onveilige invoer.

Ook de vaste e-mailschil en de verwerking van links en afbeeldingen moeten
rekening houden met ongeldige of ongewenste waarden. De gebruikte
testgegevens en screenshots mogen geen productiegegevens, wachtwoorden,
API-sleutels, tokens of andere vertrouwelijke informatie bevatten.
Gevoelige configuratiegegevens worden niet opgenomen in de repository.

De PoC introduceert geen afzonderlijk AI-model of autonome AI-functionaliteit.
De inhoud van de e-mails wordt door gebruikers aangepast en vervolgens
volgens vooraf bepaalde regels verwerkt. De Europese AI-verordening is
daarom niet de centrale focus van dit project; de relevante privacy- en
veiligheidsprincipes blijven wel van toepassing op de beoogde latere
integratie.

\section{\IfLanguageName{dutch}{Ontwerpkeuzes}{Design decisions}}%
\label{sec:ontwerpkeuzes}
Tijdens de analyse werden verschillende ontwerpkeuzes gemaakt om de
haalbaarheid van de PoC te bewaken en de oplossing voldoende eenvoudig
te houden voor de beschikbare projecttijd.

\begin{enumerate}
    \item \textbf{Een zelfstandige PoC in plaats van een rechtstreekse
    productie-integratie.}
    Een alternatief zou zijn om de editor meteen in de bestaande
    Laravel-backend en Angular-applicatie te integreren. Dat zou echter
    bijkomende afhankelijkheden creëren van de bestaande infrastructuur
    en productiearchitectuur. Door eerst een afzonderlijke mock-opstelling
    te ontwikkelen, kunnen de editor en renderingpipeline onafhankelijk
    worden getest. Het gevolg is dat een eventuele integratie later nog
    afzonderlijk moet worden ontworpen en gevalideerd.

    \item \textbf{Een vaste e-mailschil in plaats van een volledige
    layout-builder.}
    Een volledig configureerbare layout-builder zou gebruikers meer vrijheid
    bieden om de volledige structuur van een e-mail te wijzigen. Dit zou
    de complexiteit van de editor en de compatibiliteitstests echter sterk
    verhogen. Daarom blijft de e-mailschil in de PoC grotendeels vast en
    ligt de nadruk op het aanpassen van de e-mailbody.

    \item \textbf{Een visuele editor in plaats van rechtstreekse
    HTML-bewerking.}
    Bij rechtstreekse HTML-bewerking moeten gebruikers technische kennis
    hebben van HTML en de structuur van e-mails. Een WYSIWYG-editor biedt
    een toegankelijkere manier om tekst en opmaak aan te passen. De gekozen
    editor moet voldoende functionaliteit bieden voor de vereiste opmaak,
    placeholders en links, zonder onnodige complexiteit toe te voegen.

    \item \textbf{Een mock-API in plaats van de bestaande productiebackend.}
    Rechtstreeks gebruik van de Laravel-backend zou de PoC afhankelijk
    maken van bestaande bedrijfslogica, authenticatie en gegevensstructuren.
    Een afzonderlijke mock-API maakt het mogelijk de functionaliteiten
    gecontroleerd te ontwikkelen en te testen. Het nadeel is dat de
    implementatie niet zonder verdere aanpassingen naar de productieomgeving
    kan worden overgedragen.

    \item \textbf{Veilige templateverwerking met toegestane placeholders.}
    Een alternatief zou zijn om HTML en dynamische waarden zo weinig
    mogelijk te beperken. Dat vergroot echter het risico op onveilige
    invoer en onverwachte uitvoer. Daarom wordt gekozen voor HTML-opschoning,
    een whitelist van ondersteunde placeholders en veilige vervanging van
    dynamische waarden. Dit beperkt de flexibiliteit, maar verhoogt de
    voorspelbaarheid en veiligheid van de gegenereerde e-mails.

    \item \textbf{Validatie met geautomatiseerde tests en echte
    e-mailclients.}
    Alleen de HTML-output technisch controleren volstaat niet om te
    garanderen dat een e-mail in alle clients correct wordt weergegeven.
    Daarom worden geautomatiseerde tests gecombineerd met praktijktests
    in beschikbare e-mailclients. Dit levert meer inzicht op in
    compatibiliteitsproblemen, maar betekent ook dat niet iedere combinatie
    van client, apparaat, taal en template volledig kan worden getest
    binnen de beschikbare tijd.
\end{enumerate}

Deze ontwerpkeuzes ondersteunen het doel om binnen de beschikbare
projecttijd een werkende en testbare PoC op te leveren. Ze houden de
functionele scope beperkt, maken de belangrijkste technische risico's
expliciet en leveren tegelijk informatie op voor een mogelijke latere
integratie in Vesalius.ai.

